\SourceLeaf{271}%
\noindent follow. The constant must be set equal to zero, however, since $y$ vanishes when $x = 0$. Thus

\[
y = \alpha \,.\, \frac{\sin kx}{k}
\]

is the result. Hence

\[
f(x) = \alpha\, \frac{\sin kx}{k}
\]

and by 14),

\[
\tag*{(G)} g(x) = \cos kx.
\]

17. Substituting for $g(x)$ the expression $\cos kx$ in relations V. 1), 2), 3) and 4) gives the 10 relations

\[
\cos kc = \cos ka \,.\, \cos kb,
\]

\[
\begin{array}{cc}
\mathrm{tg}\, kb = \mathrm{tg}\, kc \,.\, \cos A, & \mathrm{tg}\, ka = \mathrm{tg}\, kc \,.\, \cos B, \\[4pt]
\mathrm{tg}\, ka = \mathrm{tg}\, A \,.\, \sin kb, & \mathrm{tg}\, kb = \mathrm{tg}\, B \,.\, \sin ka, \\[4pt]
\sin ka = \sin A \,.\, \sin kc, & \sin kb = \sin B \,.\, \sin kc, \\[4pt]
\cos A = \cos ka \,.\, \sin B, & \cos B = \cos kb \,.\, \sin A,
\end{array}
\]

\[
\cos kc = \cot A \,.\, \cot B,
\]

and these are the familiar equations of \textls{spherical trigonometry} when the radius of the sphere is set equal to $\dfrac{1}{k}$. They turn into the equations of \textls{non-Euclidean trigonometry} when we choose for $k$ a purely imaginary value $k = ik'$.

It is striking that \textsc{Gauss} denotes the constant arising in the integration of differential equation 15) by $k$. We cannot say with certainty whether he meant to suggest a connection with the measure of curvature, which he likewise calls $k$ in the Disquisitiones generales circa superficies curvas. A statement by \textsc{Minding} deserves mention here. It comes from his paper \emph{Beiträge zur Theorie der kürzesten Linien auf krummen Flächen}, which appeared in 1840 in \textsc{Crelle}'s Journal für die reine und angewandte Mathematik (vol. 20, pp. 323--327).

»That the formulas of spherical trigonometry hold between the sides and angles of a triangle formed by shortest lines on any surface of constant positive curvature follows at once when we recall that every surface of this kind can be developed onto a sphere. If the measure of curvature is negative, the same formulas hold with the change that the hyperbolic functions of the sides take the place of the trigonometric functions. If $a, b, c$ are the sides of the triangle, $A$ the angle opposite $a$, and $k$ the constant measure of curvature, whether positive or negative, it is easy to prove the following equation

\[
\cos a\BareRoot{k} = \cos b\BareRoot{k} \,.\, \cos c\BareRoot{k} + \sin b\BareRoot{k} \,.\, \sin c\BareRoot{k} \,.\, \cos A.\text{«}
\]

\textsc{Minding} had already determined the surfaces of revolution of constant negative curvature in 1839 (\textsc{Crelle}'s Journal, vol. 19). Compare the following note by \textsc{Gauss}, found in a notebook and dating from between 1823 and 1827.
